\documentclass[11pt,a4paper]{article}
\usepackage{amsmath,amssymb,amsthm,mathtools,booktabs,microtype}
\usepackage[margin=26mm]{geometry}
\usepackage[colorlinks=true,linkcolor=blue,citecolor=blue,urlcolor=blue]{hyperref}
\usepackage{xurl}
\setlength{\emergencystretch}{3em}
\newtheorem{theorem}{Theorem}[section]
\newtheorem{lemma}[theorem]{Lemma}
\newtheorem{proposition}[theorem]{Proposition}
\newtheorem{corollary}[theorem]{Corollary}
\theoremstyle{remark}\newtheorem{remark}[theorem]{Remark}
\newcommand{\Q}{\mathbb Q}\newcommand{\Z}{\mathbb Z}\newcommand{\F}{\mathbb F}
\DeclareMathOperator{\sqf}{sqf}\DeclareMathOperator{\Cl}{Cl}
\DeclareMathOperator{\Sel}{Sel}\newcommand{\Sha}{\mathop{\amalg}}
\title{Finite-support simultaneous twists and degree-seven pole equations}
\author{ZQ Diophantine Project}
\date{23 November 2026}
\begin{document}\maketitle
\begin{abstract}
We audit the square-denominator and prescribed-squareclass arguments in
Bhargava--Burungale, arXiv:2611.16130v2, Section~2. Their subdivided-star
construction extends to every admissible divisor of one fixed finite
prime support, with no pairwise-coprimality requirement on the divisors.
One twist parameter works for the entire family, and its extra prime
factors may avoid any fixed finite set. We deduce a common-parameter
result for native pole blocks with the same input squareclass. A
rank-one Kummer obstruction explains why this theorem cannot join two
different squareclasses congruent to seven modulo eight. Separately,
a birational rearrangement lowers the square-denominator equation from
total degree nine to seven, without lowering its three-witness count.
The existing complete defining-variable bounds are not improved by
these subsystem results alone.
\end{abstract}

\section{Audit of the arithmetic input}
For $D\ne0$ put $E_D:y^2=x^3-D^2x$. Throughout this note,
$[a/p]\in\F_2$ denotes the additive Legendre symbol at an odd prime
$p\nmid a$: it is zero for a quadratic residue and one otherwise.
All squarefree representatives of positive rational numbers are positive.

Bhargava--Burungale's Lemma~2.2 gives the implication
\begin{equation}\label{eq:BB}
 (1-c)(s^4-c^2)(t^4-c^2)=z^2
 \quad\Longrightarrow\quad c\text{ has square reduced denominator}
 \qquad(c\in\Q^\times).
\end{equation}
Its valuation argument is valid at every prime, including two, and
permits zero witnesses: if $v_p(c)<0$ is odd, the valuations of
$s^4$ and $c^2$ differ, so both quartic factors have finite even
valuation, while $1-c$ has odd valuation. The product is consequently
nonzero of odd valuation.

For the converse the exact preparation is
\begin{equation}\label{eq:prep}
 0<c<1,\qquad c\text{ has square reduced denominator},\qquad
 v_2(c)\ge3,\qquad\sqf(c)\equiv7\pmod8.
\end{equation}
The last condition implies that $v_2(c)$ is even, so the displayed
lower bound is equivalent to $v_2(c)\ge4$ here. With
$j=\sqf(c)$ and $k=\sqf(1-c)$, the square-denominator condition
is what proves $\gcd(j,k)=1$ and the mutual local-square conditions.
It must not be omitted. Also $k\equiv1\pmod8$.
The case $k=1$ is harmless, either by the source's explicit solution
or by its twist construction with an empty second leaf set.

The source's descent equations are
\begin{equation}\label{eq:descent}
 Aa=D_yb,\qquad Ab=D_yb+D_za,
\end{equation}
where $A_{pq}=[q/p]$ off the diagonal, $A\mathbf1=0$,
$y_p=[-1/p]$, and $z_p=[2/p]$. The normalization by the two
rational two-torsion points separately fixes the sign of the first
Kummer coordinate and the parity at two of the second. At a prime
$p\mid D$, the translated first coordinate has nonpositive valuation,
so $1\pm D/X$ are local squares. This gives exactly
\eqref{eq:descent}; it also handles a local two-torsion representative
by its limiting Kummer value. No surjectivity from these linear
conditions to rational points is used: positive rank is established
separately.

The source's genus criterion is supplied by the older
Tian--Yuan--Zhang theorem. For squarefree $D\equiv5\pmod8$, define
$\gamma(h)=\#(2\Cl(\Q(\sqrt{-h})))\bmod2$. Then
\begin{equation}\label{eq:genus}
 \sum_{\substack{D=d_0\cdots d_r\\
 d_0\equiv7\ (8),\ d_1\equiv3\ (8)\\
 d_i\equiv1\ (8)\ (i\ge2)}}\prod_i\gamma(d_i)=1
 \quad\Longrightarrow\quad
 \operatorname{ord}_{s=1}L(E_D,s)=1.
\end{equation}
Here $r\ge1$, all factors exceed one, and the factors after $d_1$
are unordered. The decisive precise reference is the \emph{proof}
of \cite[Theorem~3.5(2), p.~760]{TYZ}: its second bracketed sum has
$d_0\equiv7$, and the total residue forces $d_1\equiv3$.
Oddness puts the genus point outside $A(K_D)^-+A[4]$.
Lemma~3.18 of that paper shows that a torsion genus point would
belong to $A[4]$, hence the point is non-torsion. The opening
assertion of Theorem~3.5 identifies its rational tensor class with
the vector defined on p.~739, which is zero if the analytic invariant
$\mathcal L(D)$ is zero. By the definition on p.~722 and the root
number $-1$, the analytic rank is exactly one. Thus no conjectural
converse to Gross--Zagier is being substituted at this step.
The source also cites \cite[Theorem~2.1, Table~2, row~5(b)]{Smith};
the argument here was checked directly against the published
Tian--Yuan--Zhang paper and does not require an independent reading
of Smith's table.


The R\'edei minor formulas also agree with the primary matrix convention
in \cite[Theorem~3.1 and (16)]{Stevenhagen}. For $h\equiv3,7\pmod8$,
quadratic reciprocity identifies that matrix with $A_h$, whose left and
right kernels both contain $\mathbf1$. For $h\equiv1\pmod8$, ordering
the ramified prime two first gives the full matrix
$\bigl(\begin{smallmatrix}0&y^t\\z&A_h\end{smallmatrix}\bigr)$.
Deleting its first row and an odd-prime column gives
$(z\mid A_h^{(v)})$. Its determinant detects maximal possible rank,
since the deleted row is the sum of the others and the original kernel
contains $(0,\mathbf1)$. Thus the displayed minor test is precisely
vanishing of the four-rank, equivalently oddness of $\#(2\Cl)$.

Finally, a one-dimensional normalized descent space together with
analytic rank one yields algebraic rank one and finite $\Sha$ by
Gross--Zagier--Kolyvagin. The exact Selmer sequence then gives
$\Sha[2]=0$, so the finite order of $\Sha$ is odd. This does not
assert the full BSD leading-coefficient formula. No defect was
found in these Section~2 deductions. This audit does not reprove
the external genus-point or Gross--Zagier--Kolyvagin theorems.

\section{One twist for every admissible divisor of a finite support}

\begin{theorem}[Finite-support simultaneous twists]\label{thm:finite}
Let $j$ be positive squarefree with $j\equiv7\pmod8$, let $J$
be its prime support, and let $K$ be a finite set of odd primes
outside $J$. Assume $[j/q]=0$ for every $q\in K$. Set
\[
 \mathcal K=\left\{k:\ k\mid\prod_{q\in K}q,\quad
 k\equiv1\pmod8,\quad[k/p]=0\text{ for every }p\in J\right\}.
\]
There are infinitely many positive squarefree $n\equiv5\pmod8$,
with $j\mid n$ and $\gcd(n,\prod_{q\in K}q)=1$, such that for
every $k\in\mathcal K$:
\begin{enumerate}
\item $E_{nk}$ has analytic and algebraic rank one;
\item its Tate--Shafarevich group is finite of odd order;
\item it has a rational point with nonzero ordinate and abscissa in
$j\Q^{\times2}$.
\end{enumerate}
The additional prime factors of $n/j$ can avoid any prescribed finite
set. The construction has exactly $2|J|+|K|$ prime factors in $n$.
In particular, any finite family of squarefree $k_i\equiv1\pmod8$
satisfying the original mutual-square conditions with $j$ may be used
simultaneously, even when the $k_i$ have arbitrary common prime factors.
\end{theorem}

\begin{proof}
Choose $p_*\in J$ and draw a star on $J$ with centre $p_*$.
Draw a second star with centre a new prime $p_0\equiv3\pmod8$
and leaves $K$. Subdivide each of the $|J|-1+|K|$ edges by a
new prime $\ell\equiv1\pmod8$. Prescribe a nonsquare symbol
between $\ell$ and each of its two neighbours, and square symbols
between $\ell$ and every other original vertex and every other
subdivision prime. Choose the subdivision primes successively,
outside the prescribed exceptional set and $J\cup K$.
The symbol at $p_0$, which has not yet been chosen, is postponed.
For $p\in J$ prescribe
\begin{equation}\label{eq:p0}
 [p_0/p]=\deg_J(p)+y_p\quad\text{in }\F_2,
\end{equation}
where the degree is before subdivision. Also prescribe
$[p_0/\ell]=0$ on the first star and $1$ on the second star.
Reciprocity is compatible because every subdivision prime is one
modulo eight. The Chinese remainder theorem and Dirichlet's theorem
give the desired primes. With all subdivision primes fixed, $p_0$
ranges over an infinite Dirichlet set. Put
\[
 n=j p_0\prod_{\ell}\ell.
\]
All factors are distinct; the asserted residue, support, avoidance,
and prime-factor count follow.

Fix an arbitrary $k\in\mathcal K$. In the descent system for
$D=nk$, a second-star leaf $q\in K$ occurs exactly when $q\mid k$;
\emph{every subdivision prime still occurs}. At a subdivision
vertex having both original neighbours present, its equations force
the two neighbours' $a$-labels to agree and their $b$-labels to agree.
At a second-star subdivision vertex whose $K$-neighbour is absent,
its equations instead force its own labels to equal those of $p_0$.
Thus the original labels are constant on each of the two stars.
Write their two common pairs as $(a_1,b_1),(a_2,b_2)$.

The equations at the original noncentral leaves uniquely determine
the remaining subdivision labels. Summing the original rows in
either star cancels the subdivision terms. The absent-leaf terms
also cancel, because their labels equal the centre's. The resulting
system has
\[
 \overline A=\begin{pmatrix}1&1\\0&0\end{pmatrix},\qquad
 \overline y=(1,1),\qquad\overline z=(0,1).
\]
Indeed, $\sum_{p\in J}[p_0/p]=1$, since the sum of the star's
degrees is even and $\sum_{p\in J}y_p=1$.
The extra first-to-second symbols sum to $[k/p]=0$ at each
$p\in J$. In the reverse direction, reciprocity changes the
$p_0$ sum to zero, while each selected $q$ has $[j/q]=0$.
The residues of $j$ and $p_0k$ are seven and three modulo eight,
respectively. This proves the displayed reduced data for every
admissible subset of $K$.

The reduced system is the line
\[
 (a_1,b_1,a_2,b_2)=(\epsilon,\epsilon,0,0),\qquad\epsilon\in\F_2.
\]
For its nonzero vector the $a$-labels at all subdivision vertices
are zero. At a first-star leaf $p$, this follows from
\eqref{eq:p0} and $[k/p]=0$; at a second-star leaf it follows
from $[j/q]=0$; at an absent leaf's subdivision vertex it follows
from the centre label zero. Therefore the nonzero $a$-vector is
exactly the characteristic vector of $J$. The full descent space
has dimension one, because all subdivision labels are uniquely
determined and the centre equations are precisely the summed
conditions already imposed.

It remains to evaluate \eqref{eq:genus}. The determinant form of
R\'edei's formula used in \cite[(2.5)]{BB} is
\[
 \gamma(h)=
 \begin{cases}
 \det A_h^{(v,v)},&h\equiv3,7\pmod8,\\
 \det\bigl(A_h^{(v)}\mid z_h\bigr),&h\equiv1\pmod8,
 \end{cases}
\]
where $(v,v)$ deletes row and column $v$, $(v)$ deletes only
column $v$, and an empty determinant is one.
If a subdivision prime lies in a factor containing neither of its
neighbours, that factor's determinant is zero (also for a singleton
subdivision prime). If it lies in a factor with exactly one neighbour,
expansion after deleting that neighbour's column removes the subdivision
prime without changing $\gamma$.

Order the subdivided edges with both original endpoints present.
In any nonzero summand whose original endpoints are separated, take
the first such edge and move its subdivision prime to the other endpoint's
factor. The preceding deletion rule preserves the term. This operation
is an involution without fixed points; it preserves the original-vertex
partition and all residue classes, and it works with unordered
residue-one factors. A subdivision vertex with an absent $K$-neighbour
must stay with $p_0$ in every nonzero term.

Consequently only the factorization into the two complete remaining
stars survives. Their residues are seven and three. Neither can join
the other, since the product has residue five, and there is no factor
made solely of subdivision primes. For either surviving factor,
delete its centre. Each present leaf and its subdivision vertex then
form a pair whose row and column expansions contribute one. After
these pairs are removed, absent leaves of the second star leave an
identity matrix. Both determinants are therefore one. This proves
that the genus sum is odd for every $k\in\mathcal K$.

Now \eqref{eq:genus} gives analytic rank one. The rational and Selmer
Kummer images inject into the one-dimensional normalized descent
space, whose unique nonzero vector has first class $j$. Positive rank
ensures that a nonzero rational class exists, so that class is realized.
The rank and $\Sha$ assertions follow exactly as in Section~1.
Since the choices were uniform in $k$ and only $p_0$ varies, all
conclusions are simultaneous.
\end{proof}

\begin{remark}[Scope relative to earlier project work]
The earlier project finite-family theorem on the row-5(a) route uses
an input class congruent to one modulo eight and pairwise-coprime
auxiliary twists; its change-of-base argument covers some overlaps.
Theorem~\ref{thm:finite} treats the class-seven route and permits all
overlaps. It neither changes the earlier theorem's hypotheses nor
claims that class-one and class-seven inputs are interchangeable.
\end{remark}

\section{Parameter sharing and its exact limitation}
The existing native pole radicand is
\begin{equation}\label{eq:native}
 P_d(v,s)=d(d+v^2)^2s^4-d(d-v^2)^2(1-d)^2.
\end{equation}
The two-coordinate existence and finite-guard density theorem for
one prepared input was already proved in the project's v56 note.
It is not new here.

A \emph{retained parameter} means a freely chosen rational coordinate
before finite-algebra elimination, as in \cite{Common}. A square root
of a specified rational function is a finite choice, but its rationality
must still be imposed. Counting retained parameters therefore does not
by itself count the existential variables in a defining polynomial.

\begin{corollary}[Common input squareclass]\label{cor:sharing}
Let $d_1,\ldots,d_m$ satisfy \eqref{eq:prep}, and suppose they
have a common squarefree representative $j$. Let $h$ be the number
of their distinct nontrivial complement squareclasses.
Their native pole blocks admit a common construction with $h+1$
retained rational parameters, together with finite quadratic choices.
The available tuples of those retained parameters are Zariski dense.
Extra guards can be imposed whenever their polynomials remain
nonzero after the sharing substitutions.
\end{corollary}
\begin{proof}
Write $d_i=j r_i^2$ and $1-d_i=k_i h_i^2$, with
$r_i,h_i\in\Q^\times$. Preparation supplies the hypotheses of
Theorem~\ref{thm:finite} for every $k_i$. Choose its common $n$.
Choose a rational non-torsion point
$R_0=(j a_0^2,y_0)$ on $E_n$, and independently choose
$R_k=(j a_k^2,y_k)$ on $E_{nk}$ for each distinct nontrivial $k$.
Their odd multiples keep the same abscissa squareclass and supply
infinitely many choices. Set
\[
 e=\frac{n}{j a_0^2},\quad
 p=\frac{y_0}{j a_0^3},\quad v=\frac{p}{1+e},\quad
 S_k=\frac{a_k}{a_0}.
\]
Then $0<e<1$, $p^2=j(1-e^2)$, and
$e=(j-v^2)/(j+v^2)$. For the trivial complement put $S_1=1$.
Scaling $R_k$ by $(x,y)\mapsto(h_i^2x,h_i^3y)$ gives the
second point needed for the native pole block at $d_i$.
The exact substitutions are
\[
 v_i=r_i v,\qquad s_i=h_iS_{k_i},
\]
and the polynomial identity is
\begin{equation}\label{eq:sharingidentity}
 P_{d_i}(v_i,s_i)=r_i^6h_i^4j
 \bigl((j+v^2)^2S_{k_i}^4-k_i^2(j-v^2)^2\bigr).
\end{equation}
For every nontrivial $k$, the curve equation gives
\[
 \left(\frac{(j+v^2)y_k}{j a_0^2a_k}\right)^2
 =j\bigl((j+v^2)^2S_k^4-k^2(j-v^2)^2\bigr).
\]
Thus the radicands in \eqref{eq:sharingidentity} are rational squares.
For $k_i=1$, the rightmost radicand is $4j^2v^2$,
so no second free parameter is needed. Repeated complement classes
use the same $S_k$.

Varying $R_0$ supplies infinitely many $v$; for each fixed choice,
each independent $R_k$ supplies infinitely many $S_k$. Successive
coefficient comparison proves Zariski density in $\mathbb A^{h+1}$.
For a uniform coefficient construction one need not use the
squarefree integers as variable coefficients. Retain $v_*=r_1v$,
and, for a representative index $a$ of each nontrivial complement
class, retain $s_a=h_aS_{k_a}$. For an index $i$ in that class,
adjoin finite rational square roots
\[
 \alpha_i^2=d_i/d_1,\qquad
 \beta_i^2=(1-d_i)/(1-d_a),
\]
and set $v_i=\alpha_i v_*$ and $s_i=\beta_i s_a$.
The finite equation for that class is
\[
 Y_a^2=d_1(d_1+v_*^2)^2s_a^4
       -d_1(d_1-v_*^2)^2(1-d_a)^2,
\]
and the original terminal ordinate is
$Y_i=\alpha_i^3\beta_i^2Y_a$.
For a trivial complement adjoin $s_i^2=1-d_i$ and take
$Y_i=2d_iv_i(1-d_i)$.
All divisions are defined at the prepared rational input.
There are finitely many index-groupings for fixed $m$.
These formulas use only existing input functions, retained coordinates
and finite square choices; $j,k_i,n$ occur only in the existence proof.
\end{proof}

The finite choices in this corollary are not cost-free rational
witnesses. A complete construction must impose their rationality,
check the hypotheses of its finite-algebra encoder, and account for
any new generators, guards and separating coordinates. In particular,
$1-d_a$ can vanish at another algebraic character of a larger source
even when the selected prepared value satisfies $0<d_a<1$.
One must justify the required localization or handle that stratum;
density of the retained parameters cannot make this fixed zero
denominator invertible. Only guards that remain nonzero after the
sharing substitutions can be avoided. Thus $h+1$ is the retained
parameter count of this pole subsystem, and is not automatically
subtractable from a complete route's defining-variable count.

The parameter-sharing pattern was already present for the class-one
route in \cite{Common}; the new content here is the class-seven
finite-family input without an overlap restriction.

\begin{remark}[All complements square]
If $h=0$ and zero terminal ordinates are allowed, each block also has
$v_i=0$, $s_i=\sqrt{1-d_i}$ and $Y_i=0$. Thus this special locus
needs no retained free pole coordinate, only finite square choices.
This specialization has no parameter-density assertion and may violate
extra guards, especially a guard requiring a nonzero terminal ordinate.
The $h+1$ construction above remains useful when generic guard avoidance
is needed.
\end{remark}

The existing complete routes do not prove that their distinct native
pole inputs have one common squareclass. In the eleven-variable route
there is already only one native block. Thus this corollary establishes
no smaller complete defining-variable count.

The following obstruction was already recorded in the project's v82
common-level applicability note; its proof is included to make the
limitation of the new theorem explicit.

\begin{proposition}[Rank-one obstruction to mixing input classes]
Let $n>0$ be squarefree with $n\equiv5\pmod8$, and assume
$\operatorname{rank}E_n(\Q)=1$. If $E_n(\Q)$ has a non-torsion
point with abscissa in $j\Q^{\times2}$ for positive squarefree
$j\equiv7\pmod8$, then $j$ is the only positive squarefree
representative congruent to seven modulo eight in its abscissa
squareclass image. In particular two distinct such $j$ cannot be
combined on the same rank-one $E_n$.
\end{proposition}
\begin{proof}
The abscissa Kummer homomorphism $\delta_x$ has
$\delta_x(O)=1$, $\delta_x((0,0))=-1$, and
$\delta_x((n,0))=n$. It factors through $E_n(\Q)/2E_n(\Q)$,
which has order eight. The classes $-1,n,j$ are independent:
the positive torsion classes are $1,n$, of residues one and five,
whereas $j$ has residue seven. Hence the image is exactly
$\langle-1,n,j\rangle$. Its positive representatives are those of
$1,n,j,nj$, with residues $1,5,7,3$ modulo eight. Removing a common
odd square does not change the residue. Only $j$ has residue seven.
\end{proof}

The simultaneous theorem supplies a common \emph{twist parameter}.
It does not generally supply equal abscissas on its different curves,
nor does it prove a prescribed value of the actual elliptic detector
trace. The project's prepared counterexamples on $E_{5235}$,
$E_{3615}$ and $E_{791}$ therefore remain relevant. No inference
from local solubility to rational solubility is made here.

\section{A degree-seven square-denominator equation}
\begin{proposition}\label{prop:seven}
For $c\in\Q^\times$, the existence of $s,t,r\in\Q$ satisfying
\begin{equation}\label{eq:seven}
 (s^2+c)(t^2-c)=(1-c)r^2(s^2-c)(t^2+c)
\end{equation}
implies that $c$ has square reduced denominator. For every $c$
satisfying \eqref{eq:prep}, solutions exist. This equation has total
degree seven, twelve monomials and coefficient height one.
It retains three rational witnesses.
\end{proposition}
\begin{proof}
Suppose $v_p(c)=m<0$ is odd. For any rational $s$,
\[
 v_p(s^2+c)=v_p(s^2-c)=\min(2v_p(s),m),
\]
including $s=0$, and similarly for $t$. All four factors are nonzero.
If $r=0$, \eqref{eq:seven} is impossible. Otherwise the right side
has valuation equal to the left side plus $m+2v_p(r)$, a nonzero
odd integer. This is a contradiction. The argument includes $p=2$.

On the prepared locus, $c$ is a positive nonsquare and $c\ne1$.
Thus every factor used in the following divisions is nonzero.
Given \eqref{eq:BB}, set
\[
 r=\frac{z}{(1-c)(s^2-c)(t^2+c)}.
\]
It satisfies \eqref{eq:seven}. Conversely any solution of
\eqref{eq:seven} gives \eqref{eq:BB} by
$z=(1-c)r(s^2-c)(t^2+c)$. The exact residual identity is
\begin{align*}
 &(1-c)(s^4-c^2)(t^4-c^2)
   -\bigl((1-c)r(s^2-c)(t^2+c)\bigr)^2\\
 &\quad=(1-c)(s^2-c)(t^2+c)
 \bigl((s^2+c)(t^2-c)-(1-c)r^2(s^2-c)(t^2+c)\bigr).
\end{align*}
Existence follows from the audited converse in Section~1. Expansion
gives the stated degree, number of monomials and height.
\end{proof}

For example, the prepared value $c=112/225$ has
\[
 s=\frac{28}{15},\qquad t=\frac{644}{825},\qquad r=\frac{15}{29}.
\]
It corresponds to $z=657707008/765703125$ in \eqref{eq:BB}.
The shortcut $s=t$ or $s=-t$ cannot work at this $c$: it would
make $1-c=113/225$ a square, since $s^4-c^2\ne0$.
Thus the birational degree reduction does not collapse the two
independent quartic coordinates.

Equation~\eqref{eq:seven} is quadratic in $r$ with nonconstant
leading coefficient. Turning it into a monic quadratic by clearing
its leading coefficient recovers a degree-nine radicand. It must not
be advertised as an automatic improvement to the existing monic
finite-algebra compiler or to its variable counts.

\section{Exact checks and limits of the certificate}
The accompanying standard Python/SymPy checker verifies the
polynomial residual identity and the displayed rational witness.
It constructs a concrete overlapping family with
\[
 j=7,\quad K=\{113,137,193\},\quad
 n=7\cdot19\cdot449\cdot953\cdot1009=57422493709.
\]
For every one of the eight divisors $k$ of $113\cdot137\cdot193$,
it computes the normalized descent matrix over $\F_2$, checks its
one-dimensional kernel and prescribed first vector, and computes
the full unordered genus sum, obtaining one. A second example uses
$j=3\cdot5\cdot17=255$ and $K=\{409\}$.
It also checks 289 deterministic formal graph fibres, including
arbitrary reciprocity-compatible symbols and absent leaves.
Eleven small R\'edei parity computations are independently compared
with class numbers obtained by enumerating reduced primitive binary
quadratic forms. These checks support the explicit algebra and graph
implementation; the proof of the infinite-family theorem is the
argument in Section~2, using the external genus-point criterion.
The checker does not numerically compute the analytic ranks or
search for rational points on the large example curves.

\begin{thebibliography}{9}
\bibitem{BB} M. Bhargava and A. Burungale,
\emph{A Diophantine definition of $\Z$ in $\Q$},
arXiv:2611.16130v2, 19 November 2026, Section~2.
\url{https://arxiv.org/html/2611.16130v2}.
\bibitem{TYZ} Y. Tian, X. Yuan and S.-W. Zhang,
\emph{Genus periods, genus points and congruent number problem},
Asian J. Math. \textbf{21} (2017), 721--774, especially
pp.~722, 739--740, 754 and 760.
\url{https://web.math.princeton.edu/~shouwu/publications/cn.pdf}.
\bibitem{Smith} A. Smith,
\emph{The congruent numbers have positive natural density},
arXiv:1603.08479v2 (2016).
\url{https://arxiv.org/abs/1603.08479}.
\bibitem{Stevenhagen} P. Stevenhagen,
\emph{Redei reciprocity, governing fields and negative Pell},
Math. Proc. Cambridge Philos. Soc. \textbf{172} (2022), 627--654,
Theorem~3.1 and equation~(16).
\url{https://doi.org/10.1017/S0305004121000335}.
\bibitem{Native} ZQ Diophantine Project,
\emph{A two-coordinate pole test of coefficient height four},
17 November 2026, v56,
\nolinkurl{papers/pole_parameters/intrinsic_pole_20261117_v56.tex}.
\bibitem{Common} ZQ Diophantine Project,
\emph{Common twists for several pole inputs}, 14 November 2026, v30,
\nolinkurl{papers/pole_parameters/common_twists_20261114_v30.tex}.
\end{thebibliography}
\end{document}
